\section{Introduction}
\label{sec:intro}
% \addcontentsline{toc}{section}{Introduction}

In theorem provers and dependently-typed programming languages, users often have to prove to a type checker that two terms are equivalent. These proofs often involve repetitive application of the axioms defining the underlying theory. For instance, a proof of the classic identity $$\forall a, b \in \R, (a + b)^2 = a^2 + 2ab + b^2$$ invokes 10 semiring axioms: associativity of the addition three times, commutativity of the multiplication once, left distributivity of times over plus three times, left neutrality of the multiplication twice and evaluation once.

In order to free users from needing to construct such proofs, most dependently-typed languages include simplifiers for common algebraic structures. Among these structures are semirings and rings. A semiring consists of a set with two binary operations, typically called addition and multiplication, such that the semiring is a commutative monoid with respect to addition, a monoid with respect to multiplication and multiplication distributes over addition. The set of natural numbers $\N$ equipped with the usual operations form a semiring. A ring is a semiring that is an abelian group with respect to addition, i.e., each element must have an additive inverse. The set of integers $\Z$ equipped with the usual operations form a ring. Solvers for semirings and rings already exist and are implemented in most proof assistants. The state-of-the-art in terms of ring solvers is Rocq's current implementation of its \ring tactic \parencite{coq_ring}. This is also the implementation chosen by Agda \parencite{agda_solver}. The technique used is called reflection, and while \citeauthor{coq_ring} unify the implementation for semirings and rings, it is limited to those theories. For instance, one would need to implement a new solver from the ground up to deal with involutive rings.

The goal of this internship is to find a way to combine existing solvers for component theories to obtain a solver for distributive combinations of theories. Doing so would let us implement a semiring solver from a monoid solver and a commutative monoid solver. Other than letting us implement solvers more easily, this also opens the door to a wide variety of other semiring solvers, considering combinations of varieties of semigroups, monoids and groups.

This internship is part of a wider project called \textbf{fr}ee \textbf{ex}tensions, abbreviated \textsc{Frex}. Introduced in \textcite{frex_origin}, and built on universal algebra, it was first conceived as a way to represent partially-static structures, where some parts of a data structure are known at compile-time (static) and other parts are known only at run-time (dynamic). \textsc{Frex} relies on \textit{normal forms}, i.e., canonical representatives using the algebraic laws while also evaluating concrete elements. The first application of \textsc{Frex} was to optimise the code generated by multi-stage programming (MSP), a form of metaprogramming where parts of source code are annotated and become available as symbolic expressions. 
% Such staged code can be manipulated within the same program before it is translated at either run-time or compile time. This warrants the use of partially-static structures.
\textsc{Frex} can also be specialised to fully dynamic structures, considering free algebras (abbreviated fral) instead of free extensions. With its use of normal forms, \textsc{Frex} represents terms modulo semantics. This leads to so-called `representation theorems' which can be used for algebraic simplification thanks to universal algebra. The use of normal forms in the theory may also be useful for search and synthesis, by enumerating candidates over a smaller search-space.

Later, \textcite{idris_frex} formalised these representation theorems constructively and with respect to setoids, alongside other universal algebra concepts. This let them implement an extensible dependently-typed library for algebraic simplifiers based on fral and frex representation theorems.

\textsc{Frex} can be extended modularly. For instance, \parencite{involutive_frex} describes a modular construction of an involutive algebra frex out of the frex of the underlying algebra. This lets them expand on already constructed free extensions, and combined with the work of \textcite{idris_frex}, this gives solvers for involutive structures for free. For instance, implementing a simple monoid solver using frex yields an involutive monoid solver. This motivates the search for a modular construction for solvers for distributive combination of theories using \textsc{Frex}.

Reformulated using \textsc{Frex} concepts, the goal of this internship is to find a way to combine free algebras and free extensions for component theories to obtain free algebras/extensions for distributive combinations of theories.

\subsection{Outline and contributions}

In \secref{sec:preliminaries}, we present the relevant mathematical concepts behind \textsc{Frex}, how they are used in practice to simplify equations and also some concepts of category theory used to state the main result.

In \secref{sec:construction}, we present the theoretical construction of free algebras for the distributive combination of theories. Ohad Kammar had stated an incomplete result (\secref{sec:initial_result}), but it turns out it contradicted known results (\secref{sec:failure}). I combed the (incomplete) proof for errors, strengthened the hypotheses to state a new corrected result and proved it (\secref{sec:corrected_result}). We started extending the result from free algebras to free extensions, but the construction is incomplete (\secref{sec:frex_attempt}).

In \secref{sec:implementation}, we briefly present the implementation of the distributive combination of free algebras in Idris 2, what remains to be implemented, as well as a full performance evaluation (\secref{sec:perf-eval}). To do so, I also had to implement some additional solvers for the component theories (\secref{sec:solver_implem}).

In \secref{sec:comparison_main}, we compare different aspects of \textsc{Frex} with the state-of-the-art in terms of semiring solvers: Rocq's reflection-based \ring tactic \parencite{agda_solver}. \secref{sec:reflection} and \secref{sec:almost-ring} give insights into how the \ring tactic works, while \secref{sec:soundness-completeness} explains the soundness and completeness guarantees of both libraries. \secref{sec:comparison} compares the capabilities and applications of \textsc{Frex} and the \ring tactic.

The appendices contain some additional definitions and proofs (\secref{sec:additional_fral})
% , some additional context in which this internship took place (\secref{sec:context})
, the Idris 2 source code for the whole project (\secref{sec:source_code}) and the full proof of Theorem~\ref{thm:correct} (\secref{sec:full_proof}).
% and a link to my research notes (\secref{sec:notes}).